\section*{历史注记（第一、二、三章）}
\phantomsection
\addcontentsline{toc}{section}{历史注记（第一、二、三章）}

（注意：罗马数字指本注记末尾所附的参考文献。）

1604 年，正处于科学事业鼎盛时期的伽利略相信，自己已经证明：在速度随所经距离成比例增长的直线运动中，运动定律恰好就是他在重物下落中发现的那个定律\((x = ct^{2})\)（III, v. X, p.~115-116）。在 1695 至 1700 年间，每月在莱比锡出版的《Acta Eruditorum》没有一卷不刊登莱布尼茨、伯努利兄弟、洛必达侯爵的论文，这些论文以本质上仍是我们今天使用的记号，处理微分学、积分学和变分学中最多种多样的问题。因此，正是在这差不多整整一个世纪的间隔之内，锻造出了无穷小演算，或者按英国人最后的叫法，那门出类拔萃的“演算”（calculus）；而将近三个世纪的持续使用，也还没有完全磨钝这件无与伦比的工具。

希腊人既不曾拥有、也不曾想象过任何类似的东西。毫无疑问，倘若他们知道巴比伦人的那种代数演算，他们也会拒绝使用它，而他们的《几何学》有一部分很可能只是这种演算的转录；因为希腊人最辉煌的数学创造力恰恰出现在几何学发明的领域之中，那正是他们处理那些在我们看来属于积分学的问题的方法。欧多克索斯在处理锥体与棱锥体的体积时给出了这种方法最早的范例，欧几里得又以大体忠实的方式把它们传给了我们（(I)，第 XII 卷，命题 7 和 10）。但尤其是，阿基米德的几乎全部工作都献给了这些问题（(II) 与 (II bis)）；而且，由于一个奇特的机缘，他的大部分著作今天仍能以原文读到——用的是那精心结撰、音调铿锵的多利亚方言——直至最近才被重新发现的那一篇，其中他阐述了引导他得到其最美成果之一的“启发式”程序（(II), v. II, p.~425-507）。因为欧多克索斯的“穷竭法”的弱点之一正在于此：尽管它是一种无可指摘的证明方法（在约定了若干公设之后），它却不是一种发现的方法；它的运用必然以预先知道所要证明的结果为前提；因此，正如阿基米德所说，“在欧多克索斯最早找到证明的那些结果中，关于锥体与棱锥体的那些……，有相当一部分要追溯到德谟克利特，是他最先不加证明地陈述了它们”（loc. cit., p.~430）。这一情形使得对阿基米德著作的细致分析格外困难；说实在的，这种分析似乎还没有任何一位现代史学家着手去做过；因为我们并不知道，他对把他所处理的各种问题联系在一起的那些亲缘关系（即我们今天会表述为“同一个积分以各种几何面貌出现在许多地方”的那种关系）意识到什么程度，也不知道他可能赋予这些关系何种重要性。例如，考察下列问题，第一个由欧多克索斯解决，其余由阿基米德解决：棱锥体的体积、抛物线弓形的面积、三角形的重心、阿基米德螺线的面积（用极坐标即 \(\rho = c\omega\)）；它们全都依赖于积分\(\int x^2 dx\)，而且，丝毫不背离穷竭法的精神，可以把它们全都归结为形如\(\sum_{n} an^2\)的“黎曼和”的求值。阿基米德处理螺线（(II), v. II, p.~1-121）时正是这样做的，他使用了一个引理，其内容相当于写出

\[
N ^ {3} <   3 \sum_ {n = 1} ^ {N} n ^ {2} = N ^ {3} + N ^ {2} + \sum_ {n = 1} ^ {N} n <   (N + 1) ^ {3}.
\]

至于三角形的重心，他（用穷竭法，借助分解成平行条带的做法）证明了重心位于每条中线上，因而位于这些中线的公共交点上（(II), v. II, p.~261-315）。对抛物线他给出了三种做法：第一种是启发式的，只打算“给结果以某种可信性”，它借助一个静力学论证把问题归结为三角形的重心，在论证过程中他毫不犹豫地把抛物线弓形看作无穷多条平行于轴的线段之和（(II), v. II, p.~435-439）；另一种方法依据类似的原则，但以穷竭法完全严格地写出（(II), v. II, p.~261-315）；最后一种证明极为巧妙，但适用范围较小，它利用抛物线的特殊性质，把所求面积表为一个几何级数之和。没有任何迹象表明这些问题与棱锥体体积之间有联系；甚至有这样的说法（(II), v. II, p.~8）：关于螺线的问题与关于球以及旋转抛物面的另外一些问题“毫无共同之处”，而阿基米德在同一篇引言里恰好谈论过后两者，其中之一（抛物面体积）可归结为积分\(\int x dx\)。

从这些例子可以看出，除了一些特殊的技巧之外，穷竭法的原理如下：通过分解成“黎曼和”，得到所考察量的上界与下界，再把这些界与对该量所陈述的表达式直接比较，或者与一个已经解决的类似问题的相应界相比较。这种比较（由于还没有负数，必然要分两部分进行）以一句套语开场：“若不然，它就或者较大或者较小；倘属可能，设它较大，等等；倘属可能，设它较小，等等”；由此得出了“归谬法”或“化归于不可能”（“\(\dot{\alpha}\pi\alpha\gamma\omega\gamma\eta \varepsilon\dot{\iota}\sigma \dot{\alpha}\delta\acute{\upsilon}\nu\alpha\tau o\nu\)”）这一名称，那是 17\(^{th}\) 世纪的学者们给予这种方法的。阿基米德（(II), v. II, p.~62-76）对螺线切线的确定正是以类似的形式陈述的；这是一个孤立的结果，也是我们所能引证的“微分学”的古代源泉中唯一的一个，此外就只有相对容易的圆锥曲线切线的确定，以及若干极大极小问题。的确，就“积分”而言，不仅体积理论，而且静力学和流体静力学，都为希腊数学家提供了广阔无垠的研究领域；可是由于缺少运动学，他们几乎没有什么机会认真地接触微分。诚然，阿基米德给出了他的螺线的运动学定义；而且，既然我们无从知道他怎么会得到其切线的知识，我们就有权问，他是否已有运动合成的概念。但倘若如此，他难道不会把这么强有力的方法应用到同类的问题上去吗？更合理的推测是：他必定使用了某种启发式的取极限程序，这也许是他所知道的圆锥曲线的结果提示给他的；这些结果当然本质上更简单，因为直线与圆锥曲线的交点可以作出，从而可以确定这些点重合的条件。至于切线的定义，切线被设想为这样一条直线：在曲线上某点的一个邻域内，曲线完全位于它的一侧；它的存在性是假定的，而且还假定每条曲线都由凸弧组成；在这些条件下，要证明一条直线与一条曲线相切，就必须证明某些不等式，而这当然是以最完备的严格性来做的。

从严格性的观点看，阿基米德的方法无可指摘；甚至在 17\(^{th}\) 世纪，当最一丝不苟的数学家想把一个被认为特别微妙的结果置于不容置疑的地位时，他们给出的仍是“归谬法”的证明（(XI a) 与 (XII a)）。至于这些方法的多产性，阿基米德的全部著述就是充分的见证。但是，若要有理由在其中看出一门“积分学”，就必须在多种多样的几何外观之下，展示出某种按其底下所含“积分”的性质对问题进行分类的雏形。我们将看到，在 17\(^{th}\) 世纪，对这样一种分类的探求逐渐成为几何学家们主要的关注之一；如果在阿基米德那里找不到一点痕迹，这难道不正好是一个标志，表明这类思辨在他看来会显得过于“抽象”，而恰恰相反，他在每一个场合都刻意尽可能贴近他所研究的图形的特殊性质吗？而我们难道不能得出这样的结论：这部辉煌的著作——按照创立者们的说法，积分学完全取自于它——在某种意义上竟是积分学的反面？

此外，在数学中任凭发现与证明之间出现一道鸿沟，是不可能不付代价的。在顺利的时代，数学家不必放弃严格性，只需把他的想法几乎照构思时的样子写下来；有时他甚至可以指望凭借对通行的语言和记号作一次恰当的调整而做到这一点。但更常见的情形是，他不得不在两者之间作出取舍：一边是虽可能富有成果却并不正确的表述方式，一边是那些只有把他的思想扭曲了才能表达出来、且须付出令人疲惫的努力的正确方法。两条路都各有风险。希腊人走了第二条路，而他们的数学在其最辉煌的绽放之后几乎立刻就令人惊讶地停滞了，其原因或许正应当到这里来寻找，其分量甚至超过罗马征服的毁灭性作用。有人提出过一个并非没有道理的猜想：阿基米德和阿波罗尼奥斯的后继者们的口头讲授中想必已包含许多新结果，只是他们认为不值得强迫自己付出那种按公认规范发表所必需的非凡努力。无论如何，使 17\(^{th}\) 世纪数学家止步的并不是这些顾忌：面对纷至沓来的新问题，他们在阿基米德的著作中孜孜不倦地寻找克服这些问题的手段。

希腊文学与哲学的古典名著已由阿尔杜斯·马努提乌斯及其后继者在意大利全部刊印，而且几乎都在 1520 年以前；而阿基米德的第一个希腊文与拉丁文对照版本\(^{4}\)直到 1544 年才由赫尔瓦吉乌斯在巴塞尔印行，此前并没有任何拉丁文译本先行；而且，当时的数学家（他们正专心于各自的代数研究）远没有立刻感受到这一影响，一直要等到伽利略和开普勒——两人都更是天文学家和物理学家而非数学家——这种影响才变得显著。从这时起，直到大约 1670 年不曾间断，无穷小演算创立者们的著作中出现得最频繁的名字，莫过于阿基米德。许多人翻译他、为他撰写评注；从费马到巴罗，人人都不加区别地引证他；人人都声称在那里同时找到了典范和灵感之源。

诚然，我们将会看到，这些声明并非全都可以完全按字面理解；这里就遇到了妨碍正确解释这些著作的困难之一。历史学家还必须考虑到当时科学界的组织状况：17\(^{th}\) 世纪初它还非常不完善，而到该世纪末，由于学术团体和科学期刊的创立，由于大学的巩固和发展，它最终已与我们今天所了解的十分相似。在 1665 年之前没有任何期刊，数学家除了写信，就是印一本书——多半自费，或者倘能找到一位梅塞纳斯式的赞助人，便由他出资——此外别无办法让世人了解自己的工作。能承担这类工作的编辑和印刷商很少，有时也不可靠。在这样一次出版所必然带来的长期拖延和无数变故之后，作者还常常不得不面对无休止的论战：论战由并非总是心怀善意的对手挑起，有时以出人意料地尖刻的口气进行；因为，在无穷小演算的基本原理普遍不确定的情况下，任何人要在对手的论证中找出薄弱之处、至少是含混而可争议之处，都不困难。可以理解，在这种情形下，许多珍视安宁的学者满足于把他们的方法和结果告知少数几位经过挑选的朋友。有一些人，尤其是若干科学爱好者，例如巴黎的梅森以及后来伦敦的柯林斯，与各国保持着大量通信，并把其中的摘录四处转告，其间也并非不掺入一些他们自己发明的错误。数学家们手中握有一些“方法”，但由于缺乏概念和一般的定义，他们既不能把这些方法写成定理的形式，甚至也无法把它们确切地表述出来，于是只得在大量特例上检验它们，并且相信没有比向同行们发出挑战更能较量彼此实力的了，有时还伴随以用密码语言发表自己的结果。好学的年轻人四处游学，其程度也许甚于今日；一位学者的思想有时更多地是靠其学生的游历而不是靠他本人的著作传播开来，而这又提供了误解的另一个来源。最后，同样的问题自然会出现在许多数学家面前，其中有些人非常杰出，却只是不完整地了解旁人的结果，于是优先权之争不绝如缕，而且往往还要加上剽窃的指控。

因此，历史学家必须到这一时期学者们的书信和私人文稿中去寻找自己的文献，其分量几乎不亚于、甚至超过他们的正式出版物。然而，例如惠更斯的书信已得到保存并成为典范性出版的对象（XVI），莱布尼茨的却至今只以有缺陷而零碎的方式刊行，还有许多别的书信已无可挽回地散失。至少，以手稿分析为依据的最新研究，已经以一种似乎无可辩驳的方式澄清了一个几乎被派系之争掩盖的事实：每当这一时期的一位大数学家谈到自己的工作、自己思想的发展演变、什么影响了他而什么没有时，他都是诚实而真诚地、完完全全出于诚意地做的\(^{5}\)；因此，这些我们掌握着相当数量的珍贵证言可以放心使用，历史学家在此无须把自己变成一位预审法官。至于其余，大多数优先权问题根本没有什么意义。的确，莱布尼茨采用 dx 这个记号表示“微分”时，并不知道牛顿用\(\dot{x}\)表示“流数”已有十年左右：可是，即使他知道，又当如何呢？举一个更有教益的例子：定理\(\log x = \int dx/x\)的作者是谁？它的日期又是什么？我们刚才所写的这个公式应归功于莱布尼茨，因为两项都是用他的记号写出的。莱布尼茨本人以及沃利斯把它归于格雷古瓦·德·圣樊尚。而后者在其《几何著作》（Opus Geometricum）（IX）中（该书出版于 1647 年，但据他自己说，写成远在此前）只证明了下述命题的等价形式：设\(f(a, b)\)表示双曲线弓形\(a \leqslant x \leqslant b\)、\(0 \leqslant y \leqslant A/x\)的面积，则关系\(b'/a' = (b/a)^n\)蕴涵\(f(a', b') = n f(a, b)\)；他的学生兼评注者萨拉萨几乎随即补充了评注\(^{6}\)：面积\(f(a, b)\)由此可以“代替对数”。如果他对这件事只说了这么多，如果格雷古瓦本人对此一言不发，这难道不是因为，对这一时期的大多数数学家来说，对数还是一些“计算的辅助工具”，尚未归化入籍，无权在数学中被引证吗？诚然，托里拆利在 1644 年的一封信（VII bis）中谈到他对一条曲线的研究，这条曲线我们会写成\(y = ae^{-cx}\)、\(x \geqslant 0\)，并补充说纳皮尔（他对之多有赞誉）“只追求实用的算术”，而他自己却“从中引出了一种几何的思辨”；他还留下一份显然为发表而准备的关于这条曲线的手稿，但它直到 1900 年才刊行（(VII), v. I, p.~335-347）。此外，笛卡尔在大约 1639 年就因“德博纳问题”遇到过同一条曲线，并且在描述它时没有提到对数（(X), v. II, p.~514-517）。无论如何，J. 格雷戈里在 1667 年给出了一条用（十进）对数计算双曲线弓形面积的法则，没有引证任何人（(XVII a)，转载于 (XVI bis), p.~407-462）：这既表明他已在理论上知道双曲线求积与对数之间的联系，又表明他已在数值上知道“自然”对数与“十进”对数之间的联系。惠更斯直接质疑格雷戈里结果之新颖性的主张（XVI a），是否仅针对最后这一点呢？这对我们来说，正如对他们的同时代人一样不清楚；无论如何，同时代人都清楚地感到，对数与双曲线求积之间存在联系这件事是早已知道的事情，只是他们除了书信中的暗示、乃至格雷古瓦·德·圣樊尚的书之外，引不出别的依据。1668 年，布龙克尔给出了 log 2 和\(\log(5/4)\)的级数（附有一个通过与几何级数比较而作的细致的收敛证明）（XIV），他把这些级数作为双曲线相应弓形的表达式提出，并补充说他所得到的数值与 2 和\(5/4\)的“对数成相同的比”。但同一年，随着墨卡托（XIII）（或者更确切地说，随着沃利斯立即对墨卡托工作所作的阐述（XV bis）），语言改变了：既然双曲线的弓形与对数成比例，又既然大家知道对数由其特征性质确定时只差一个常数因子，那么就没有什么妨碍把双曲线的弓形看作对数，称之为“自然的”（与“人工的”或“十进的”对数相对）或“双曲的”对数；这最后一步一旦迈出（墨卡托给出的\(\log(1 + x)\)的级数对此有所贡献），定理\(\log x = \int dx / x\)便得到了——只差记号——或者不如说，它甚至成了定义。由此能得出什么结论呢？无非是：这一发现是经过几乎察觉不到的过渡而完成的，而在这件事上的优先权之争，极像小提琴与长号为一支交响曲中某一动机出现的准确时刻而起的争吵。而且说实话，尽管同时代的其他数学创造——费马的数论、牛顿的运动学——都带有强烈的个人印记，17\(^{th}\) 世纪无穷小演算的发展让人想起的，却是一部交响曲逐渐而必然的展开：其中“时代精神”（Zeitgeist）既是作曲家又是指挥，执掌着指挥棒；每个人都以自己独有的音色奏出自己的声部，但谁也不是他所奏出的那些主题的主人，这些主题几乎难以分解地被精妙的对位缠绕在一起。因此，这段历史只能以主题分析的形式来写；我们在此仅满足于一个概略的草图，并不奢求分毫不差的精确\(^{7}\)。下面至少就是在粗略考察下呈现出的主要主题：

\begin{enumerate}
\def\labelenumi{\Alph{enumi})}
\tightlist
\item
  数学严格性的主题，与无穷小、不可分量或微分的主题相对照。我们已经看到，这两者在阿基米德那里都占有重要地位：前者贯穿于他的全部著作，后者则只出现在他唯一的一部《方法》（Method）论著中，而 17\(^{th}\) 世纪并不知道这部论著，因此，如果说它被传递了下来而不是被重新发明，那就只能是通过哲学的传统。此外，无穷小的原理以两种不同的形式出现，视其涉及的是“微分”还是“积分”而定。关于这一点，首先假设要计算一块平面面积：人们将用无穷多条等距的平行线把它分成无穷多条无穷小的平行条带；而每条这样的条带都是矩形（尽管用两条相距有限距离的平行线所得到的有限条带并不是矩形）。类似地，旋转体将被垂直于轴的平面分解成无穷多个具有同一无穷小高度的圆柱\(^{8}\)；在用共点的直线把一块面积分解成三角形时，或者在讨论一段曲线弧的长度时把它当作无穷多边的多边形来处理时，等等，人们也会使用类似的说话方式。可以肯定，那少数几位牢固掌握了阿基米德方法的数学家，如费马、帕斯卡、惠更斯、巴罗，在每一个具体场合都不会感到用严格的证明来代替这种语言有什么困难；因而他们也常指出，这种语言不过是一种速记。费马说：“按阿基米德的方式给出证明并不难；……，只要事先一次性地交代清楚，就可以避免不断重复……”（(XI), v. I, p.~257）；帕斯卡同样说：“因此这两种方法之一与另一种并无差别，只是说法不同而已”（(XII b), p.~352）\(^{9}\)；还有巴罗，以他那种讥诮的简洁说道：“longior discursus apagogicus adhiberi possit, sed quorsum?”（本来可以再添一段归谬式的论述，但那有什么用呢？）（(XVIII), p.~251）。费马似乎不愿提出他无法如此论证的东西，因而也就只能以暗示的方式或以“方法”的形式陈述任何一般性结果；巴罗尽管十分谨慎，却稍微不那么拘泥。至于他们同时代人中的大多数，至少可以说，严格性并不是他们主要关心的事情，而阿基米德的名字在大多数时候只是一顶帐篷，用来遮盖那些无疑价值高昂、但阿基米德
\end{enumerate}

肯定不会为之承担责任的货物。在涉及微分的时候就更是如此。如果把待求长的曲线当作一个具有无穷多条边的多边形，那就是把曲线的一段“无穷小”弧当作一段“无穷小”线段——或者是弦，或者是其存在性已被假定的切线上的线段；或者反过来，人们考虑的是一个“无穷小”的时间区间，在这区间内（当人们只涉及速度时）运动“是”均匀的；更大胆的是笛卡尔，为了确定摆线的切线——这不在他的普遍法则范围之内——他把相互滚动的曲线都当作多边形，从而推断“在无穷小之中”运动可以看作绕接触点的旋转（(X), v. II, p.~307-338）。在这里，费马把他的切线法则以及极大极小法则建立在这样的无穷小考虑之上，却仍能在每个具体场合为它们给出论证（(XI b)；另见 (XI), v. II, 各处，特别是 p.~154-162，以及《著作补编》（Supplément aux Œuvres，Gauthier-Villars, 1922），p.~72-86）；巴罗则按古人的方式，从单调性和凸性这样简单的假设出发，对他的许多定理给出了精确的证明。但那时已不再是往旧皮囊里灌新酒的时候了。正如我们今天所知道的，这一切之中逐渐成形的是极限概念；即使人们能从帕斯卡、牛顿以及其他人的著作中找出一些似乎非常接近我们现代定义的陈述，也只需把它们放回原来的语境，就能看出那些阻碍严格表述的无法逾越的障碍。当 18\(^{th}\) 世纪起一些讲究清晰的数学家想在他们那堆乱七八糟的财富中理出些条理时，他们前辈著作中的这类提示对他们极为宝贵；例如，当达朗贝尔解释说微分中除了极限概念之外别无他物，并对它下了精确的定义时（XXVI），可以相信他受到了牛顿关于“消失量的最初比与最终比”的论述（XX）的指引。但是，就 17\(^{th}\) 世纪而言，必须非常明确地指出：只有当牛顿和莱布尼茨背弃过去、暂时同意不在严格证明中、而是在结果的丰富多产和首尾一致之中去寻求新方法的辩护时，通向现代分析的道路才算打开。

\begin{enumerate}
\def\labelenumi{\Alph{enumi})}
\setcounter{enumi}{1}
\tightlist
\item
  运动学。前面已经看到，阿基米德已经给出了他的螺线的运动学定义；而在中世纪发展出了一种（但在缺乏相反证据的情况下，并无无穷小考虑的）关于量随时间变化及其图像表示的初步理论，它也许可以追溯到巴比伦的天文学。但对 17\(^{th}\) 世纪的数学至关重要的一点是：微分问题从一开始出现，就不仅关系到切线，而且关系到速度。然而，伽利略（(III) 与 (III bis)）在研究重物下落中速度的规律时（他此前已通过斜面实验得到了距离定律\(x = at^{2}\)），并不是通过微分来进行的：他对速度作出各种假设，先是\(v = dx/dt = cx\)（(III), v. VIII, p.~203），后来又是\(v = ct\)（同上，p.~208），然后试图通过对速度随时间变化的图像进行推理——方式相当含混——来重新得出距离定律；笛卡尔（1618 年）对定律 v = ct 的论证与之类似，但他是作为一个真正的数学家、以不可分量语言所能容许的最大清晰度来做的\(^{10}\)（(X), v. X, p.~75-78）；在这两人那里，速度的图像（此处是一条直线）起着主要作用，而人们可以问，他们究竟在多大程度上意识到所经距离与时间轴和速度曲线之间所围面积成比例；不过很难就这一点断言什么，尽管笛卡尔的语言似乎表明他知道上述事实（某些历史学家想把它追溯到中世纪\(^{11}\)），而伽利略对它则没有明确的暗示。巴罗在 1670 年明确地陈述了它（(XVIII), p.~171）；对当时任何人来说这也许都不算新，巴罗也没有把它当作新东西提出；但是，无论对这一结果还是对任何别的结果，都不应该试图把日期定得过分精确。至于伽利略也考虑过的假设\(v = cx\)，他（同上）只是满足于证明它是站不住的（或者用现代的语言说，方程\(dx/dt = cx\)没有既\(\neq 0\)又在\(t = 0\)处为零的解），其论证相当含混，费马后来（(XI), v. II, p.~267-276）不厌其烦地把它展开（这差不多等于说，既然\(2x\)与\(x\)同为解，那么\(x \neq 0\)就与物理上显而易见的解的唯一性相违背）。而正是这同一条定律\(dx/dt = cx\)，在 1614 年被纳皮尔用来引进他的对数；他对对数给出了一种运动学的定义（IV），用我们的记号可写成：若两条直线上各有一个动点分别按定律\(dx/dt = a\)、\(dy/dt = -ay/r\)、\(x_0 = 0\)、\(y_0 = r\)运动，则称\(x\)为\(y\)的“对数”（用现代记号即\(x = r \log(r/y)\)）。我们已经看到，方程\(dy/dx = c/x\)的解曲线于 1639 年出现在笛卡尔那里，他对它作了运动学的描述（(X), v. II, p.~514-517）；的确，他把一切非代数曲线相当轻蔑地归为“机械的”，并声称要把它们逐出几何学；但所幸这一禁忌——很久以后莱布尼茨还认为必须对它激烈抗议——既没有被他的同时代人遵守，也没有被笛卡尔本人遵守。摆线和对数螺线出现了，并且得到了热切的研究，其研究极大地促进了几何方法与运动学方法的相互渗透。运动合成的原理，更确切地说速度合成的原理，是伽利略在其晚年杰作、1638 年的《谈话录》（Discorsi）（(III), v. VIII, p.~268-313）中所阐述的抛体运动理论的基础；这一理论因此隐含地包含了抛物线切线的一种新的确定方法；如果说伽利略没有明确指出这一点，托里拆利（(VII), v. III, p.~103-159）则相反，他坚持了这一点，并且正是以这一原理为基础，建立了一种对那些能给出运动学定义的曲线确定切线的一般方法。诚然，他在这件事上被罗贝瓦尔（VIII \(a\)）抢先了几年，后者说他正是通过研究摆线而想到这种方法的；也正是摆线切线这同一个问题，给了费马显示其微分方法之威力的机会（(XI b), p.~162-165）；而笛卡尔由于在这里无法应用他的代数方法，便为此专门发明了瞬时转动中心（(X), v. II, p.~307-338）。
\end{enumerate}

但是，随着无穷小演算的发展，运动学不再是一门独立的科学。人们越来越清楚地看到，尽管有笛卡尔，代数曲线和代数函数从无穷小演算的“局部”观点来看，并没有任何东西使它们有别于其他远为一般的曲线和函数；具有运动学定义的函数和曲线与其他函数和曲线一样，同样可以应用相同的方法；而变量“时间”不过是一个参数，它的时间性纯粹是语言上的事情。因此，在惠更斯那里，即使处理的是力学问题，占统治地位的仍是几何（XVI b）；莱布尼茨在他的演算中也没有赋予时间任何特殊的地位。与此相反，巴罗从各种量作为一个被设想为“时间”的独立普适变量的函数而同时变化出发，构造出一种具有几何倾向的无穷小演算的基础。这个想法——想必是他在试图重新找出运动合成方法时想到的，他对这种方法的存在性只是耳闻——在他的《几何讲义》（Lectiones Geometricae）（XVIII）的前三讲中得到了非常清晰而有力的详细阐述；例如，他在那里仔细地证明：如果一个动点在两条互相垂直的轴 AY、AZ 上的投影是两个动点，其中一个以常速度 a 运动，另一个以随时间增加的速度 v 运动，那么轨迹有一条斜率等于 v/a 的切线，并且朝 Z 增大的方向凹。在《几何讲义》的其余部分，他把这些思想发挥得非常深远；尽管他执意从头到尾以一种尽可能几何而少代数的形式来表述它们，人们仍可以像雅各布·伯努利那样（(XXIII), v. I, p.~431 和 453），在其中看到牛顿和莱布尼茨无穷小演算中相当大一部分的等价物。恰恰是同样的思想成了牛顿的出发点（(XIX c) 与 (XX)）：他的“流量”（fluentes）是各种量，是一个仅仅作为普适参数的“时间”的函数，而“流数”（fluxions）就是它们对“时间”的导数；牛顿允许自己视需要更换参数，这种做法在巴罗的工作中同样存在，只是使用得不如他灵活\(^{12}\)。流数的语言由牛顿采用，并凭借他的权威强加给了下一世纪的英国数学家，因此，就我们所处理的时期而言，它代表了其真正作用业已终结的运动学方法的最后一个据点。

\begin{enumerate}
\def\labelenumi{\Alph{enumi})}
\setcounter{enumi}{2}
\tightlist
\item
  代数几何。这是一个寄生性的主题，与我们的论题无关，它之所以闯进来，是因为笛卡尔以他那系统化的态度，声称要把代数曲线作为几何学的唯一对象（(X), v. VI, p.~390）；因此，他为确定切线所给出的方法是代数几何的方法，而不像费马那样是来自微分学的方法。古人遗留下来的关于直线与圆锥曲线相交的结果、笛卡尔本人关于两条圆锥曲线相交的思考，以及可归结为它们的问题，本来会极其自然地把他引到以两个交点重合作为相切判据的想法：我们今天知道，在代数几何中这是正确的判据，而且其普遍性如此之大，以至于它与极限概念以及“基域”的性质无关。笛卡尔最初以一种不太方便的方式应用它：设法让所研究的曲线与一个圆心在 Ox 上的圆的两个交点在给定点重合（(X), v. VI, p.~413-424）；他的学生范·斯霍滕和赫德把圆换成直线，得到了曲线
\end{enumerate}

\[
\mathrm{F} (x, y) = 0
\]

的切线斜率，其形式为\(-F_{x}^{\prime}/F_{y}^{\prime}\)，其中“导出多项式”\(F_{x}^{\prime}\)、\(F_{y}^{\prime}\)由它们的形式构成规则定义（(X bis), v. I, p.~147-344 和 (XXII), p.~234-237）；德·斯吕塞在大约同一时间也独立得到了这一结果（(XXII), p.~232-234）。当然，我们在此指出的那些截然分明的区别——唯有它们才使笛卡尔与费马之间的争论有意义——在 17\(^{th}\) 世纪数学家的头脑中是不可能以任何方式存在的：我们提到它们，只是为了澄清与我们所关心的历史相关的最奇特的事件之一，并指出代数方法随即就完全销声匿迹了，它们暂时被微分方法吸收了。

\begin{enumerate}
\def\labelenumi{\Alph{enumi})}
\setcounter{enumi}{3}
\tightlist
\item
  问题的分类。我们已经看到，这一主题在阿基米德的著作中似乎是缺失的，而且对一个问题是直接解决还是把它归结为一个已经处理过的问题，在他看来是无关紧要的。在 17\(^{th}\) 世纪，微分问题最初以三种不同的面貌出现：速度、切线、极大与极小。至于最后这类问题，开普勒（V）注意到（这一观察在奥雷姆\(^{13}\)那里已经出现，甚至巴比伦的天文学家也没有放过它）：函数的变化在极大值附近特别缓慢。费马在 1630 年以前（(XI bis)；参见 (XI), v. II, p.~71）就这些问题开创了他的无穷小方法，用现代语言说，这相当于考察泰勒展开的前两项（常数项和一阶项），并写出第二项在极值处为零；随后他把方法推广到切线的确定，甚至应用于拐点的寻求。如果把上面关于运动学所说的话考虑在内，就可以看出，这三类关于一阶导数的问题的统一是相当早就完成的。至于涉及二阶导数的问题，它们出现得较晚，尤其见于惠更斯关于曲线渐屈线的工作（发表于 1673 年他的《摆钟论》（Horologium Oscillatorium）（XVI b）中）；这时牛顿凭借他的流数，已经掌握了解决这类问题所需的全部分析工具；而尽管惠更斯在这上面倾注了全部几何才干（后来的微分几何在其初创时期从中受益），在我们所处理的时期内，这些问题除了让新分析得以宣告其工具的威力之外，并无别的用处。
\end{enumerate}

至于积分，它在希腊人那里表现为面积、体积、矩的计算，表现为圆的周长和球冠面积的计算；17\(^{th}\) 世纪又增加了曲线的求长、旋转曲面面积的计算，以及（随着惠更斯关于复摆的工作（XVI b））转动惯量的计算。首先需要认识的是所有这些问题之间的联系。就面积和体积而言，第一步巨大的跨越是由卡瓦列里在他的《不可分量的几何》（VI a）中迈出的。他在书中陈述并声称为证明了大致如下的原理：如果两个平面面积具有这样的性质——每条平行于给定方向的直线截它们所得的线段成恒定的比——那么这两个面积也成同样的比；对于被平行于固定平面的平面所截、所得诸截面度量成恒定比的体积，也陈述了类似的原理。这些原理很可能是由欧几里得（或者毋宁说欧多克索斯）关于同高棱锥体积之比的定理提示给卡瓦列里的，而且在以一般方式陈述它们之前，他想必先在取自阿基米德著作的大量例子上验证过它们的有效性。他使用一种语言来“论证”这些原理——关于这种语言的合法性，人们看到他 1621 年还在一封信里征询伽利略的意见，尽管 1622 年起他已毫不犹豫地使用它了（(III), v. XIII, p.~81 和 86）——这种语言的要点如下。例如考察两块面积，一块是\(0 \leqslant x \leqslant a\)、\(0 \leqslant y \leqslant f(x)\)，另一块是\(0 \leqslant x \leqslant a\)、\(0 \leqslant y \leqslant g(x)\)；纵坐标之和\(\sum_{k=0}^{n-1} f(ka/n)\)与\(\sum_{k=0}^{n-1} g(ka/n)\)之比，当\(n\)充分大时可以任意接近两块面积之比，甚至当 f 和 g 单调时，用穷竭法证明这一点也不难；卡瓦列里取了极限，令\(n = \infty\)，并谈到第一条曲线“所有纵坐标之和”，它与第二条曲线相应的和之比严格等于两块面积之比；体积情形相同；这种语言后来被普遍采用，甚至连费马这样对它所涵盖的精确事实有最清楚认识的作者也不例外。的确，后来许多数学家，如罗贝瓦尔（VIII a）和帕斯卡（XII b），宁愿把对其取“和”的曲线的纵坐标不看作卡瓦列里那样的线段，而看作一些具有同一无穷小宽度的矩形；从严格性的观点看这算不上多大的进步（不管罗贝瓦尔怎么说），但也许能使想象力不至于太容易出轨。无论如何，由于涉及的只是比，曲线\(y = f(x)\)的“所有纵坐标之和”，简称曲线的“全部纵坐标”，归根结底——正如它也出现在帕斯卡的著作中那样——正是莱布尼茨记号\(\int y dx\)的确切等价物。

从卡瓦列里采用的语言出发，上述原理必然成立，而且立即蕴含一些我们即将用现代记号陈述的结论，这里约定\(\int f dx\)就是指 Ox 与曲线\(y = f(x)\)之间所围的面积。首先，任何平面面积，只要每条直线 x = 常数截它所得诸线段长度之和为\(f(x)\)，它就等于\(\int f dx\)；任何体积，只要每个平面 x = 常数截它所得截面的测度为\(f(x)\)，结论同样成立。其次，\(\int f dx\)线性地依赖于 f；我们有\(\int (f + g) dx = \int f dx + \int g dx\)，\(\int cf dx = c \int f dx\)。特别地，一切面积问题和体积问题都归结为求积，也就是归结为形如\(\int f dx\)的面积的计算；而且，也许更有新意也更重要的是，必须把依赖于同一求积的两个问题视为等价，并且人们拥有工具来逐例判定情况是否如此。希腊数学家从未达到（或许永远也不会同意达到）这样的“抽象”程度。于是（(VI), p.~133）卡瓦列里相当容易地“证明”了两个相似立体之比等于相似比的立方，而阿基米德对旋转二次曲面及其截体陈述这一结论时，已是在他关于这些立体的理论的末尾（(II), v. p.~258）。但要达到这一步，就必须把阿基米德式的严格性抛进海里。

这样，人们就有了——至少暂时地——按问题所归结成的求积的实际或表面困难程度来对问题分类的手段。这正是当时的代数所充当的样板：因为在代数中，在几何学中出现的代数问题中，尽管希腊人只关心解答本身，16\(^{th}\) 和 17\(^{th}\) 世纪的代数学家已开始把主要注意力转向按可用于求解的工具的性质对问题进行分类，从而预示了现代的代数扩张理论；他们不仅按问题所依赖的方程的次数对问题作了初步的分类，而且已经提出了关于可能性的困难问题：用根式解每一个方程的可能性（许多人已不再相信这一点）等等（见第 II 卷第 V 章的历史注记）；他们还致力于把给定次数的所有问题化归为一种几何类型的形式。同样，在积分学中，卡瓦列里的原理使他能够立即认出，阿基米德解决的许多问题都可归结为 n = 1, 2, 3 时的求积\(\int x^{n} dx\)；他还设计了一种巧妙的方法，对想要多少个 n 值就能对多少个 n 值完成这一求积（这种方法相当于按齐次性观察到\(\int_{0}^{2a} x^{n} dx = c_{n} a^{n+1}\)，并写出

\[
\int_ {0} ^ {2 a} x ^ {n} d x = \int_ {- a} ^ {a} (a + x) ^ {n} d x = \int_ {0} ^ {a} \left((a + x) ^ {n} + (a - x) ^ {n}\right) d x,
\]

展开后由此便得到\(c_{n}\)的一个递推关系（(VI a), p.~159 和 (VI b), p.~269-273）。但费马此前已走得更远：他先（1636 年以前）证明了当 n 为正整数时\(\int_{0}^{a} x^{n} dx = \frac{a^{n+1}}{n+1}\)（(XI), v. II, p.~83），用的是关于前 N 个整数幂和的一个公式（一种仿照阿基米德对螺线求积的做法）；随后又把这个公式推广到了一切有理数\(n \neq -1\)（(XI), v. I, p.~195-198）；对最后一个结果（1644 年已告知卡瓦列里），他直到很久以后、在读过帕斯卡关于积分的著作之后才写出证明\({}^{14}\)（XI c）。

这些结果，再加上起着变量替换和分部积分作用的几何考虑，已经使人们能够解决大量可归结为初等求积的问题。在此之外，人们首先遇到的是圆和双曲线的求积：由于这一时期人们主要同“不定积分”打交道，用现代的话说，这些问题的解分别由反三角函数和对数提供；这些函数是以几何方式给出的，而我们已看到后者是如何一步一步被引入分析的。这些求积成了大量著作的主题，作者有格雷古瓦·德·圣樊尚（IX）、惠更斯（(XVI c) 与 (XVI d)）、沃利斯（XV a）、格雷戈里（XVII a）；前者相信自己完成了圆的求积，后者相信自己证明了\(\pi\)的超越性；他们发展了对圆函数和对数函数逐次逼近的方法，有的偏于理论，有的面向数值计算，这些方法不久便随着牛顿（(XIX a) 与 (XIX b)）、墨卡托（XIII）、J. 格雷戈里（XVII bis）、再到莱布尼茨（XXII）而发展为级数展开的一般方法。在每一种情形下，人们都逐渐确信这些求积的“不可能性”，也就是说，确信它们所定义的函数的非代数性；与此同时，人们习以为常地把一个已被归结为某种“不可能”求积的问题，看作在其本性允许的范围内已经解决。例如摆线的问题就属于这种情形，它由三角函数解出；抛物线的求长也属此类，它被归结为双曲线的求积。

求长问题——我们刚才提到了其中两个最著名的——具有特殊的重要性，因为它们构成了一道自然的几何桥梁，连接着以其为前提的微分和它们从属的积分；旋转曲面面积的问题也可以与它们联系起来。古人只处理过圆和球的情形。这些问题在 17\(^{th}\) 世纪出现得非常晚；看来是椭圆求长（椭圆被认为是仅次于圆的最简单曲线）当时无法逾越的困难使人望而却步。运动学的方法使人们对这些问题有了一些抓手，使罗贝瓦尔（VIII b）和托里拆利（(VII), v. III, p.~103-159）在 1640 至 1645 年间得到了关于螺线弧的一些结果；但直到临近 1660 年的那些年头，这些问题才成为当务之急；摆线由雷恩于 1658 年求长（(XV), v. I, p.~533-541）；稍后，曲线\(y^{3} = ax^{2}\)由好几位作者处理（(XV), v. I, p.~551-553；(X bis), p.~517-520；(XI d)），另外还有几位（(XI), v. I, p.~199；(XVI), v. II, p.~224）把抛物线的求长归结为双曲线的求积（也就是归结为一个代数—对数函数）。最后这个例子最重要，因为它是一般原理的一个特例：曲线\(y = f(x)\)的求长不外乎\(y = \sqrt{1 + (f'(x))^2}\)的求积；而埃拉也确实是从这一原理把它推出来的。同样有趣的是追寻暮年的费马在其关于曲线\(y^3 = ax^2\)的工作（XI d）中的尝试；他把曲线\(y = f(x)\)与其弧长 s = g(x)所对应的曲线\(y = g(x)\)联系起来，并由第一条曲线的切线确定了第二条曲线的切线（用现代的语言说，他证明了它们的斜率\(f'(x)\)与\(g'(x)\)由关系式

\[
\left(g ^ {\prime} (x)\right) ^ {2} = 1 + \left(f ^ {\prime} (x)\right) ^ {2});
\]

人们觉得自己已经非常接近巴罗，只需把这个结果与埃拉的结果结合起来（格雷戈里 1668 年所做的差不多正是这件事（(XVII bis), p.~488-491)），便可得到切线与求积之间的关系；但费马只陈述了这样一点：如果两条曲线各自参照一组直角坐标轴时，横坐标相同的点处的切线斜率总相同，那么这两条曲线相等，或者换句话说，\(f'(x)\)的知识确定\(f(x)\)（相差一个常数）；而他对这一断言的论证只是一个含混的、毫无证明价值的推理。

不到十年之后，巴罗的《几何讲义》（Lectiones Geometricae）（XVIII）问世了。开篇（第 I 讲）他就原则上陈述了：在直线运动中，距离与时间轴和速度图像之间所围的面积\(\int_{0}^{t} v dt\)成比例。人们本以为他会由此、并由前面已引述的他关于确定切线的运动学方法，推出被视为切线斜率的导数与被视为面积的积分之间的联系；可是全然不是这样，他后来以纯几何的方式证明（(XVIII), Lect. X, §11, p.~243）：若两条曲线\(y = f(x)\)与\(Y = F(x)\)使得纵坐标 Y 与面积\(\int_{a}^{x} y dx\)成比例（也就是说，若\(cF(x) = \int_{a}^{x} f(x) dx\)），则\(Y = F(x)\)的切线与 Ox 相交于横坐标为 x - T 的点，其中 y/Y = c/T；而且证明从\(f(x)\)单调这一明确假设出发，是完全正确的：书中还指出，\(f(x)\)的变化方向确定了\(Y = F(x)\)凹性的方向。但可以注意的是，这条定理湮没在一大堆其他定理之中，其中许多也非常有趣；毫无准备的读者只会把它看成用求积解问题\(Y/T = f(x)/c\)的一种手段，也就是说，看成由切线信息确定曲线这一特殊问题（用我们的话说，即一类特殊的微分方程）的一个特例；更何况巴罗对它的应用首先涉及的正是同类问题（也就是可用“变量分离”求解的微分方程）。巴罗强加给自己的几何语言，在这里成了这样一个原因：使微分与积分之间那层在运动学中如此清楚的联系，多少被遮蔽了。

另一方面，把一些积分问题化归为另一些积分问题、“解决”它们、或者至少把它们归结为已分类的“不可能”问题的各种方法已经成形。分部积分最简单的几何形式，就是把由 Ox、Oy 和单调曲线\(y = f(x)\)连接 Ox 上一点\((a, 0)\)与 Oy 上一点\((0, b)\)的一段弧所围的面积写成\(\int_{0}^{a} y dx = \int_{0}^{b} x dy\)；而且它时常被不自觉地使用着。下述推广出现在帕斯卡（(XII c), 287-288）那里，但已藏得很好：设\(f(x)\)如上，令\(g(x)\)为一个\(\geqslant 0\)的函数，并设\(G(x) = \int_{0}^{x} g(x) dx\)；则有\(\int_{0}^{a} y g(x) dx = \int_{0}^{b} G(x) dy\)，他对此的证明十分巧妙：用两种方式计算立体\(0 \leqslant x \leqslant a\)、\(0 \leqslant y \leqslant f(x)\)、\(0 \leqslant z \leqslant g(x)\)的体积；特例\(g(x) = x^{n}\)、\(G(x) = \frac{x^{n+1}}{n+1}\)在帕斯卡（loc. cit., p.~289-291）和费马（(XI), v. I, p.~271）那里都起着重要作用；后者（其著作有一个意味深长的标题：《Transmutation et émendation des équations des courbes, et ses applications variées à la comparaison des espaces curvilignes entre eux et avec les espaces rectilignes \ldots》）没有证明它，无疑是因为他觉得不值得重复帕斯卡刚刚发表的东西。这些“变换术”定理——在我们看来是分部积分与变量替换的结合——在一定程度上取代了变量替换，因为后者直到很久以后才被引入；这确实与当时的思维方式相悖：那时的思维还太几何、太少解析，不容许使用图形所限定的变量（即两个坐标之一，或有时是极坐标，或曲线的弧长）以外的变量。正是这样，我们在帕斯卡（XII d）那里发现了一些结果，用现代记号、令\(x = \cos t\)、\(y = \sin t\)、并对特殊的函数\(f(x)\)，它们可以写成：

\[
\int_ {0} ^ {1} f (x) d x = \int_ {0} ^ {\pi / 2} f (x) y d t,
\]

而在 J. 格雷戈里（(XVII bis), p.~489）那里，对一条曲线\(y = f(x)\)及其弧\(s\)，有\(\int y ds = \int z dx\)，其中\(z = y\sqrt{1 + y^2}\)。直到 1669 年我们才看到巴罗掌握了变量替换的一般定理（(XVIII), p.~298-299）；他的陈述一如既往是几何的，相当于下述内容：设\(x\)与\(y\)由一个单调关系联系，\(p\)是该关系在点\((x, y)\)处的斜率；那么，若函数\(f(x)\)与\(g(y)\)使得关系\(f(x)/g(y) = p\)对每一对对应的值\((x, y)\)都成立，则取在相应限之间的面积\(\int f(x) dx\)与\(\int g(y) dy\)相等；反过来，若这些面积总相等（默认\(f\)与\(g\)具有常数符号），则\(p = f(x)/g(y)\)；逆命题自然可用来把该定理应用于微分方程的求解（通过“变量分离”）。但巴罗只把这条定理放进一个附录（Lect. XII, app. III, theor. IV），在那里他指出自己先前的许多结果只是它的特例，并为发现它太晚、来不及更多地利用它而表示歉意。

于是，大约在 1670 年，形势如下：人们已经会用统一的程序处理可归结为一阶导数的问题，惠更斯也已处理了可归结为二阶导数的几何问题。人们已会把所有积分问题都归结为求积；对于分类尚未明朗的情形，已有若干几何面貌的技巧可把一种求积化为另一种，并且人们已从这个观点出发习惯了三角函数和对数函数的运用；人们已意识到微分与积分之间的联系；人们已开始处理“反切线方法”——这是当时对可归结为一阶微分方程的问题的称呼。墨卡托对级数\(\log(1+x)=-\sum_{1}^{\infty}(-x)^{n}/n\)的轰动性发现，刚刚为把级数、主要是幂级数应用于那些曾被称为“不可能”的问题开辟了全新的前景。另一方面，数学家的队伍已大为稀疏：巴罗辞去了教授席位去当布道士；除惠更斯外（他的数学著述几乎已全部完成，已得到《摆钟论》的全部主要结果，正在着手做最后的编订），只有剑桥的牛顿和阿伯丁的 J. 格雷戈里还在活跃；莱布尼茨不久也将带着新入道者的热忱加入他们。这三人——牛顿早在 1665 年、J. 格雷戈里自 1668 年墨卡托的著作出版起、莱布尼茨大约自 1673 年起——都把主要精力投入到当时的热门课题即幂级数的研究上。但是，从问题分类的观点看，新方法的主要效果似乎是抹掉了问题之间的一切区别；的确，牛顿——他更是个分析学家而非代数学家——毫不迟疑地在 1676 年告诉莱布尼茨（(XXII), p.~224），他会解所有的微分方程\({}^{15}\)；莱布尼茨则回答（(XXII), p.~248-249）说，他恰恰相反，关心的是“在承认求积的前提下”尽可能用有限形式得到解，同时也想知道是否每种求积都能像大多数已研究情形所声称的那样归结为圆和双曲线的求积；他借此提到，格雷戈里相信（正如我们今天所知，他有理由相信）椭圆和双曲线的求长不能归结为圆和双曲线的求积；莱布尼茨还问到，牛顿所用的级数方法究竟能在多大程度上对这些问题给出回答。牛顿方面（(XXII), p.~209-211）则宣称自己握有判据——他并未说明——可（显然是通过考察级数）判定某些求积（有限形式）的“可能性”，并对积分\(\int x^{a}(1+x)^{\beta}dx\)给出了一个（非常有趣的）级数例子。

人们可以看到不到十年间取得的巨大进步：分类问题在这些信件中已经用完全现代的措辞提出；如果说莱布尼茨提出的那个问题已由\(XIX^{th}\)世纪的阿贝尔积分理论解决，那么另一个问题，即能否把给定的微分方程归结为求积的问题，尽管近来有不少重要工作，仍然悬而未决。情况之所以如此，是因为牛顿和莱布尼茨各自都已把无穷小演算的基本运算归结为一个算法；只需用两人中任一人的记号写下一个求积问题或一个微分方程，它的代数结构便立即显现，从几何外衣中脱胎而出；“变换”方法同样可以用简单的解析语言写出；分类问题也得到了精确的陈述。从数学上讲，17\(^{th}\) 世纪这时已经走到了尽头。

\begin{enumerate}
\def\labelenumi{\Alph{enumi})}
\setcounter{enumi}{4}
\tightlist
\item
  内插与差分演算。这一主题（我们不把二项式系数的研究与它分开）出现得很早，并贯穿整个世纪，其原因兼有理论和实际两个方面。事实上，当时的伟大任务之一就是三角函数表、对数表和航海表的计算，地理学、航海术、理论和实用天文学、物理学以及天体力学的迅速进展使这一任务成为必需；许多最杰出的数学家，从开普勒到惠更斯和牛顿，都参与了这项工作，或是直接参与，或是通过关于最有效的逼近过程的理论研究。
\end{enumerate}

在表格的使用乃至编制中，最初的问题之一就是内插；而随着计算精度的提高，17\(^{th}\) 世纪人们已经察觉，古老的内插程序在一阶差分（表中相继数值之差）不再大体保持常数时便失去了效力；于是我们看到布里格斯例如\({}^{16}\)在对数的计算中使用了高阶差分，甚至相当高阶的差分。后来，我们看到牛顿（(XIX d) 与 (XX)，第 III 卷，引理 5）\({}^{17}\)和 J. 格雷戈里（(XVII bis), p.~119-120）各自独立地并行推进他们在内插和幂级数方面的研究；两人通过各不相同的方法，一方面得到了用多项式作内插的公式，即所谓“牛顿型”的公式，另一方面得到了二项级数（(XVII bis), p.~131；(XXII), p.~180）以及经典分析的主要幂级数展开（(XVII bis)；(XIX a and d) 与 (XXII), p.~179-192 和 203-225）；这两条研究路线相互促进，而且在牛顿心中还与无穷小演算原理的发现紧密相连，这几乎是不容置疑的。格雷戈里同牛顿一样，对实际的数值工作表现出极大的关切，体现在表格的编制与使用、级数与积分的数值计算上；特别是，尽管找不到像前面引述过的布龙克尔勋爵那种细致的收敛证明，两人都不断从级数是否适合计算这一实际观点出发谈论他们级数的收敛性。我们再次看到，牛顿为回答柯林斯出于实际目的提出的问题\({}^{18}\)，把所谓欧拉–麦克劳林求和方法应用于 N 很大时\(\sum_{p=1}^{N}\frac{1}{n+p}\)的近似计算。

人们同样很早就遇到由差分计算函数值的方法，它被用作一种实际的积分程序，甚至可以说是微分方程的积分程序。比如，赖特在 1599 年为编制航海表而必须解决一个问题，我们今天会把这个问题写成

\[
{\frac {d x}{d t}} = \sec t = {\frac {1}{\cos t}},
\]

他按一弧秒的间隔相继把\(\sec t\)的值累加起来\(^{19}\)，实际上自然而然地得到了一张\(\log \tan (\pi / 4 + t / 2)\)的数值表；而这种巧合——自最早的\(\log \tan t\)表的计算以来就被人们注意到——直到 1668 年格雷戈里积分\(\sec t\)之前，一直没有得到解释（(XVII \(c\) ) 与 (XVII bis) p.~7 和 463）。

但这些问题还有一个纯理论的、甚至是算术的方面。我们约定用\(\Delta^{r}x_{n}\)表示序列\((x_{n})_{n\in\mathbb{N}}\)的逐阶差分序列，它们借助\(\Delta x_{n}=x_{n+1}-x_{n}\)递归地定义，并

\[
\varDelta^ {r} x _ {n} = \varDelta \big (\varDelta^ {r - 1} x _ {n} \big),
\]

用\(S^{r}\)表示\(\Delta\)及其迭代的逆运算，即置\(y_{n} = Sx_{n}\)（如果\(y_{0} = 0\)、\(\Delta y_{n} = x_{n}\)），并且\(S^{r}x_{n} = S\left(S^{r-1}x_{n}\right)\)；我们有\(S^{r}x_{n} = \sum_{p=0}^{n-r} \binom{n-p-1}{r-1} x_{p}\)，特别地，若对一切 n 有\(x_{n} = 1\)，则\(Sx_{n} = n\)，而序列\(S^{2}x_{n}\)与\(S^{3}x_{n}\)就是希腊算术家早已研究过的“三角形数”和“棱锥数”的序列，并且一般地\(S^{r}x_{n} = \binom{n}{r}\)对\(n \geqslant r\)成立（而\(S^{r}x_{n} = 0\)，当 n \textless{} r）；从这一观点看，这些序列肯定早在 16\(^{th}\) 世纪就已引入；它们还自然而然地出现在组合问题之中——这些问题或凭其自身、或因与概率相联系，在 17\(^{th}\) 世纪的数学中起着相当大的作用，例如在费马和帕斯卡那里，随后在莱布尼茨那里。它们还出现在前 N 个整数的\(m^{th}\)次幂之和的表达式中；正如我们已看到的，这个和的计算乃是费马第一种方法中积分\(\int x^{m} dx\)（m 为整数）的基础（(XI), v. II, p.~83）。1655 年，沃利斯在他的《无穷算术》（Arithmetica Infinitorum）（XV a）中就是这样做的：他不知道费马的（未发表的）工作，而且据他自己说，除读过托里拆利的著作外，对不可分量方法也一无所知；的确，急于求成的沃利斯没有在细致的研究上耽搁：一旦对 m 的最初几个整数值得到了结果，他便“凭归纳”假定它对一切整数 m 都成立；对 m = 1/n（n 为整数）的情形他作了正确的过渡，然后又以一种比第一次更加粗疏的“归纳”，过渡到任意的有理数 m。但这项工作的趣味和独创性在于：他由此出发逐步上升到对“欧拉”积分 I(m, n) = \(\int_{0}^{1}(1 - x^{1/m})^{n} dx\)（当 m 和 n \textgreater{} 0时，其值为\(\Gamma(m+1)\Gamma(n+1)/\Gamma(m+n+1)\)）以及其他类似积分的研究，编制了整数 m、n 时 1/I(m, n) 的数值表——它不是别的，正是整数\(\binom{m+n}{n}\)的表——并以与今天阐述\(\varGamma\)函数理论时所用的几乎相同的方法，最终得到\(\mathrm{I}(\frac{1}{2},\frac{1}{2})=\pi/4=\left(\varGamma(\frac{3}{2})\right)^{2}\)的无穷乘积；此外，只要借助分部积分和非常简单的变量替换，并考虑 m、n 取一切实数值时的\(\mathrm{I}(m,n)\)，就不难把他的方法变得严格，而这后一点他几乎不可能想到，牛顿的分析不久便使之成为可能。无论如何，正是沃利斯对整数

\[
\binom{m + n}{n}
\]

向 m 的非整数值（更确切地说，向形如 n = p/2、其中 p 为奇整数的值）所作的“内插”，成了牛顿起步时的出发点（(XXII), p.~204-206）：它引导他先通过研究特例\((1 - x^{2})^{p/2}\)得到二项级数，继而引进对一切实数 a 的\(x^{a}\)（即如此记写的），并借助二项级数对\(x^{a}\)求导；所有这些都未曾花多大力气去求得证明乃至严格的定义；此外，一项引人注目的革新是，他由\(x^{a}\)导数的知识推出了\(\int x^{a} dx\)（\(a \neq -1\)）（(XIX a) 与 (XXII), p.~225）。至于其他方面，尽管他不久就掌握了展开幂级数的远为一般的方法，例如所谓牛顿折线法（用于代数函数）（(XXII), p.~221）和待定系数法，但他后来仍带着一种偏爱，一再回到二项级数及其推广上；例如，他似乎正是由此导出了上面提到的\(\int x^{a}(1 + x)^{\beta} dx\)的展开（(XXII), p.~209）。

然而这些思想在欧洲大陆上的演变却很不相同，而且抽象得多。帕斯卡在研究二项式系数（他由之构成并命名了“算术三角形”）及其在概率演算和差分演算中的应用时，已接近费马；当他着手处理积分时，又把这些同样的思想引入其中。与他的前辈们一样，在使用不可分量语言时，他把积分\(F(x) = \int_{0}^{x} f(x) dx\)设想为“曲线的所有纵坐标之和”

\[
\mathrm{S} \left(f \left(\frac {n}{N}\right)\right) = \sum_ {0 \leqslant p <   \mathrm{N} x} f \left(\frac {p}{N}\right),
\]

与“单位”\(N = \sum_{0 \leqslant p < N} 1\)（当\(N = \infty\)时）之比的值（(XII b), p.~352-355）（或者，当他放弃这种语言而改用正确的穷竭法语言时，即该比在 N 无限增大时的极限）。但是，由于着眼于矩的问题，他注意到：当处理等距分布的离散质量\(y_i\)时，总质量的计算就归结为上面定义的运算\(Sy_n\)，而矩的计算归结为运算\(S^2 y_n\)；于是，通过类比，他迭代运算\(\int\)，构成了他所说的“纵坐标的三角形和”，用我们的语言说即和\(N^{-2} S^2(f(n/N))\)的极限，也就是积分\(F_2(x) = \int_0^x F(x) dx\)；再迭代一次便得到“棱锥形和”\(F_3(x) = \int_0^x F_2(x) dx\)，即\(N^{-3} S^3(f(n/N))\)的极限；上下文表明，他到此为止并非因为其思想或语言缺乏独创性，而只是因为他预计只会用到这些和，其系统运用正是他大部分结果的基础，而且他立即证明了这些和的性质，用我们的记号写出即\(\mathrm{F}_{2}(x)=\int_{0}^{x}(x-u)f(u)du,\quad\mathrm{F}_{3}(x)=\frac{1}{2}\int_{0}^{x}(x-u)^{2}f(u)du\)（(XII b), p.~361-367）；这一切都没有写下一个公式，但所用语言如此透明、如此精确，使人可以像我们刚才所做的那样立即把它译成公式。在帕斯卡那里，同在他的前辈们那里一样——只是清楚得多、系统得多——自变量的选取（总取两个坐标之一，或取曲线的弧）隐含在确定积分区间等距（虽然“无限接近”）分点的约定之中；这些分点视情形或者位于 Ox 上，或者位于 Oy 上，或者位于曲线弧上，而帕斯卡注意不在此问题上留下任何含糊（(XII b), p.~368-369）。当他需要换变量时，他借助一条原理，其内容相当于说：面积\(\int f(x)dx\)总可以写成\(\mathrm{S}\big(f(x_{i})\Delta x_{i}\big)\)，不论把积分区间分割成的一些“无穷小”区间\(\Delta x_{i}\)相等与否（(XII d), p.~61-68）。

可以看到，我们已经非常接近莱布尼茨了；而他当初希望入门现代数学时就遇到了惠更斯，这可以说是幸运的偶然，惠更斯立即把帕斯卡的著作放到了他手中（(XXII), p.~407-408）；他对组合分析的思考使他有备而来，我们知道他曾对其作过深入研究，这反映在他的全部著述中。1675 年，我们看到他把上面引述的帕斯卡定理抄写成\(\mathrm{omn}(x\omega) = x.\mathrm{omn}\omega -\mathrm{omn}(\mathrm{omn}\omega)\)的形式，其中\(\mathrm{omn}\omega\)是\(\omega\)从 0 到\(x\)的积分的缩写；几天之后，莱布尼茨用\(\int \omega\)（“summa omnium \(\omega\)”的首字母）代替了它，同时引入\(d\)表示无穷小的“差分”，或者像他不久以后所说的，微分（(XXII), p.~147-167）。他把这些“差分”设想为彼此可以比较、却不能与有限量比较的量，然而他多半明确或含蓄地把自变量的微分\(dx\)取作单位（\(x\)为自变量），即\(dx = 1\)（这相当于把微分\(dy\)与导数\(dy / dx\)等同起来），因而起初在积分记号中把它省去，于是积分写作\(\int y\)而不是\(\int ydx\)；但他随即引入了后者，并在一旦察觉它对自变量选取的不变性之后便系统地沿用它——这种不变性使人无须时刻惦记着这一选取\({}^{20}\)；当他重新研读此前一直忽视的巴罗时，注意到巴罗如此看重的变量替换一般定理可以立即从自己的记号推出，其满意之情溢于言表（(XXII), p.~412）。此外，在这一切中他始终紧扣差分演算，他的微分演算正是由差分演算通过一个极限过程推得的，当然，要严格论证这一过程他本得大费周章；在其后的著述中，他刻意坚持这样一个事实：他的原理对两者同样适用。例如，在与约翰·伯努利的通信（(XXI), v. III, p.~156）中提及自己的早期研究时，他明确引证了帕斯卡，给出了差分演算中的一个公式——它是牛顿公式的一个特例——并由之“取极限”得到公式\(y = \sum_{1}^{\infty} (-1)^{n-1} \frac{d^{n}y}{dx^{n}} \frac{x^{n}}{n!}\)（其中 y 是在 x = 0 处为零的函数，而\(d^{n}y/dx^{n}\)是它在变量的值 x 处的各阶导数），该公式等价于伯努利刚刚通信告诉他的一个类似公式（(XXI), v. III, p.~150) 与 (XXIV), v. I, p.~125-128），后者用逐次分部积分证明了它。可以看到，这个公式已非常接近泰勒级数；而泰勒 1715 年为得到“他的”级数而重新发现的，也正是莱布尼茨这一从差分演算出发取极限的论证\(^{21}\)，只是他后来并未对这级数多加利用。

\begin{enumerate}
\def\labelenumi{\Alph{enumi})}
\setcounter{enumi}{5}
\tightlist
\item
  从上面描述的演变中，人们已经可以看出无穷小演算逐步代数化的过程，也就是说，把它化为一种带有统一的代数记号体系的运算演算。正如莱布尼茨曾多次极其明晰地指出的（(XXI b) p.~230-233），问题在于为新分析去做韦达为方程理论所做的事、笛卡尔为几何学所做的事。要理解这样做的必要性，只须读几页巴罗的书：任何时候都离不开一幅事先详细描述的、有时相当复杂的图形；构成其著作实质的那 100 页（第 V-XII 讲）所用的图形不少于 180 幅。
\end{enumerate}

诚然，在多种多样的几何外观尚未显现出某种统一性之前，谈不上什么代数化。不过，格雷古瓦·德·圣樊尚（IX）已经（以“ductus plani in planum”之名）引入了一种合成律，它相当于系统地使用积分\(\int_{a}^{b}f(x)g(x)dx\)，即把它看作立体\(a \leqslant x \leqslant b\)、\(0 \leqslant y \leqslant f(x)\)、\(0 \leqslant z \leqslant g(x)\)的体积；但他远没有推出帕斯卡后来（如前所见）从研究同一立体中推出的那些结论。1655 年的沃利斯和 1658 年的帕斯卡各自以自己的方式锻造了具有代数特征的语言，他们在其中不写任何公式，却写下了只要领会其机制便可立即转录为积分学公式的陈述。帕斯卡的语言尤为清晰精确；尽管人们不理解他为什么拒绝使用代数记号——不仅是笛卡尔的，甚至连韦达的也不用——却只能叹服他所完成的这一杰作，而若无驾驭语言的本领，这是不可能做到的。

然而几年过去，一切都变了。牛顿第一个想到，把当时无穷小分析中所有几何性质的运算，换成单一的分析运算即微分，再加上反问题的求解；这种运算当然可以借助幂级数方法极其便利地完成。借用他的语言，我们已经看到，凭借一个普适“时间”参数的设想，他把作为该参数之函数的各种变化的量称为“流量”，把它们对时间的导数称为“流数”。他似乎并不特别看重记号，他的崇拜者后来甚至把没有系统的记号吹捧成一个优点；然而为自用起见，他不久便养成了用一点表示流数的习惯，于是 dx/dt 记作\(\dot{x}\)，\(d^{2}x/dt^{2}\)记作\(\ddot{x}\)，等等。至于积分，牛顿似乎同巴罗一样，从未把它看作一种运算，而只看作一个问题（已知流数求流量，解\(\dot{x}=f(t)\)）；因此积分没有名称，似乎也没有标准的记号（只是有时用一个小方框\(\left|\overline{f(t)}\right|\)或\(\Box f(t)\)表示\(\int f(t)dt\)）。这是因为他对一个并非唯一确定、只确定到相差一个可加常数的东西，不愿给予名称和记号吗？由于缺乏文字材料，我们只能提出这个问题。

牛顿如何地依靠经验、力求正确、即便在最大胆之处也审慎有加，莱布尼茨就如何地系统化、喜于推广、是敢于冒险的革新者，有时还颇为自负。从青年时代起，他头脑中就怀有一个“特征术”（characteristic）或普适符号语言的想法：它之于人类思想的全体，恰如代数记号之于代数；其中每个名称或符号都是所指事物全部性质的钥匙，而且只要使用得正确，就必然推理得正确。人们尽可把这样一个计划斥为空想；然而其作者恰恰就是不久之后认出并分离出无穷小演算诸基本概念、并赋予它那差不多已是最终形式记号的人，这并非偶然。上面我们已经见证了这套记号的诞生，并看到似乎自觉负有使命的莱布尼茨如何精心地逐步修改它们，以保证他所追求的简洁、尤其是不变性（XXI a 和 b）。这里需要指出的是，他从引入\(\int\)和 d、积分和微分之时起，就对它们怀有明确的概念：把它们看作互逆的算子（那时他对牛顿的思想还一无所知）。诚然，这样行事使他无法摆脱不定积分所固有的歧义——那是他体系的弱点，他机智地绕了过去，他的后继者们也是如此。但从新符号一出现起就引人注目的是，莱布尼茨立即着手表述它们的使用规则：他问是否有\(d(xy) = dx dy\)（(XXII), p.~16-166），自己作了否定的回答，随后逐步得到正确的法则（XXI a），后来又把它推广成著名的\(d^{n}(xy)\)公式（(XXI), v. III, p.~175）。当然，在莱布尼茨如此摸索之时，牛顿知道 z = xy 蕴涵\(\dot{z} = \dot{x}y + x\dot{y}\)已经十年了；但他从不觉得有必要把它说出来，因为在他看来，这不过是他对流量间关系\(F(x, y, z) = 0\)求导法则的一个不值得命名的特例。相反，莱布尼茨主要关心的，并不是让他的方法服务于这类具体问题的求解，甚至也不是从严格而无懈可击的原理把它们推演出来，而首先是建立一种以少数简单规则的形式操作为基础的算法。正是本着这种精神，他借助括号改进了代数记号，逐步采用\(\log x\)或 lx 来表示对数\({}^{22}\)，并大力提倡“指数演算”，即对\(a^{x}\)、\(x^{x}\)、\(x^{y}\)这类指数为变量的指数式作系统的考察。尤其，牛顿只在每个具体情形确属必要时才引入高阶流数，莱布尼茨却很早就致力于通过 d 与\(\int\)的迭代创造一种“运算演算”；他逐渐意识到数的乘法与其演算中算子的复合之间的类比，便以可喜的胆识采用指数记号来写 d 的迭代，即用\(d^{n}\)表示\(n^{th}\)次迭代（(XXII), p.~595 和 601\({}^{23}\)，以及 (XXI), v. V, p.~221 和 378），甚至用\(d^{-1}\)、\(d^{-n}\)表示\(\int\)及其迭代（(XXI), v. III, p.~167）；他甚至还尝试给\(d^{\alpha}\)（\(\alpha\)取任意实数值）赋予意义（(XXI), v. II, p.~301-302，以及 v. III, p.~228）。

这并不是说莱布尼茨不也关心其演算的应用；他深知应用是试金石（惠更斯就常这样向他重复（(XXII), p.~599））；但他缺乏深入其中的耐心，尤其热衷于寻找机会表述新的普遍法则。就是这样，1686 年（XXI c）他处理了曲线的曲率和密切圆问题，到 1692 年（XXI d）得到了平面曲线接触的一般原理\({}^{24}\)；1692 年（XXI e）和 1694 年（XXI f），他奠定了包络理论的基础；1702 年和 1703 年，他与约翰·伯努利各自用分解为简单元素的方法实现了有理分式的积分，但起初只是形式地进行，并未充分认识分母中出现复线性因子时所伴随的那些情形（XXI, g 和 h）。又是这样，1697 年 8 月的一天，他在马车中思索变分学的问题时，想到了在\(\int\)号下对参数求导数的法则，并立即满怀热情地把它寄给了伯努利（(XXI), v. III, p.~449-454）。而当他达到这一步时，其演算的基本原理早已确立，其使用也已开始传播：无穷小分析的代数化已是既成事实。

\begin{enumerate}
\def\labelenumi{\Alph{enumi})}
\setcounter{enumi}{6}
\tightlist
\item
  函数概念在 17\(^{th}\) 世纪以多种方式被引入和澄清。全部运动学都建立在一个直观的、某种意义上是经验性的观念之上，即随时间变化的量、也就是时间的函数的观念；我们已经看到，人们如何由此达到作为参数的函数，它出现在巴罗那里，也以“流量”之名出现在牛顿那里。“任意曲线”的概念经常出现，却很少被精确化；很可能，人们想到它时往往采取的是一种运动学的、至少是经验性的形式，并且不认为一条曲线要能充当论证的对象就必须容许几何的或解析的刻画；特别在涉及积分时（例如在卡瓦列里、帕斯卡和巴罗那里）情形就是如此——其原因我们今天能够更好地理解了；最后这位（巴罗）在讨论由 x = ct、\(y = f(t)\)定义的曲线时，在 dy/dt 递增这一假设之下，甚至明确说过，dy/dt“按某一规律的规则方式、甚至以不规则方式”增大，都“无关紧要”（(XVIII), p.~191）——用我们的话说，就是它容不容许一个解析定义都无所谓。遗憾的是，这个清晰而富有成果的观念——经过适当的澄清之后，它将在 19\(^{th}\) 世纪重新出现——当时却敌不过笛卡尔造成的混乱：笛卡尔一方面把一切不容许精确解析定义的曲线逐出“几何学”，另一方面又把这样一种定义中容许使用的生成手续局限于代数运算。的确，在最后这一点上，他的大多数同时代人并没有追随他；各种超越运算——对数、指数、三角函数、求积、微分方程的求解、取极限、级数求和——逐渐地、而且往往是经由非常迂回曲折的途径确立了起来，尽管在每一场合都不容易指出迈出这一步的确切时刻；况且，迈出第一步之后往往又跟着倒退一步。例如就对数而言，对数曲线（视坐标轴的选取为\(y = a^{x}\)或\(y = \log x\)）的出现、对数螺线的出现、双曲线的求积、\(\log(1+x)\)的级数展开、乃至\(\log x\)或 lx 这一记号的采用，都应视为重要的阶段。至于三角函数——它们在某种意义上可以一直追溯到古代——有趣的是可以观察到：正弦曲线最初并不是作为由方程\(y = \sin x\)定义的曲线出现的，而是在罗贝瓦尔那里（(VIII a), p.~63-65）作为“摆线的伴随曲线”出现的（也就是作为曲线
\end{enumerate}

\[
y = \mathrm{R} \left(1 - \cos {\frac {x}{\mathrm{R}}}\right),
\]

的一个特例），也就是说，作为一条其定义由摆线之定义导出的辅助曲线出现的。要遇到解析式的一般概念，我们必须求之于 J. 格雷戈里，他在 1667 年（(XVI bis), p.~413）把它定义为这样一个量：它由其他的量经过一系列代数运算“或任何别的可以想象的运算”而得到；他在其序言（(XVI bis), p.~408-409）中试图使这一概念精确化，解释说必须在代数的五种运算\(^{25}\)之外再添上第六种运算，而归根结底，这第六种运算不是别的，正是取极限。但这些有趣的想法很快就被遗忘了，淹没在格雷戈里本人、牛顿以及其他人所发现的级数展开的洪流之中；这后一方法异乎寻常的成功，造成了容许解析定义的函数与可展开为幂级数的函数之间经久不消的混淆。

至于莱布尼茨，他似乎持守笛卡尔的观点，只是明显地加进了求积，又含蓄地加进了当时分析所熟悉的其他运算：幂级数求和、微分方程求解。同样，约翰·伯努利在想要考虑 x 的任意函数时，把它说成“以任意方式由 x 和常数形成的量”（(XXI), v. II, p.~150），有时还具体说明所说的是“以代数方式或超越方式”形成的量（(XXI), v. II, p.~324）；1698 年，他与莱布尼茨约定把这样一个量称为“x 的函数”（(XXI), v. III, p.~507-510 和 p.~525-526）\(^{26}\)。莱布尼茨此前已经引入了“常量”、“变量”、“参数”这些词，并借包络问题精炼了依赖于一个或多个参数的曲线族的概念（XXI e）。记号问题也在他与约翰·伯努利的通信中得到了澄清：后者乐于用 X 或\(\xi\)表示 x 的任意函数（(XXI), v. III, p.~531）；莱布尼茨表示赞同，但也提议用\(x^{-|1|}, x^{-|2|}\)，即我们会写成\(f_{1}(x), f_{2}(x)\)的那种记号；而对于 x 的函数 z 的导数 dz/dx，他提出用 dx 这个记号（以区别于作为微分的 dz），伯努利则写作\(\Delta z\)（(XXI), v. III, p.~537 和 526）。

就这样，随着这个世纪的过去，英雄时代结束了。新演算连同它的概念和记号，以莱布尼茨赋予它的形式确立了下来。最早的弟子雅各布·伯努利和约翰·伯努利与宗师在发现上争先，在他为之指出路径的丰饶疆域中纵情遨游。第一部微分学和积分学论著是约翰·伯努利在 1691 年和 1692 年间为一位侯爵写的\(^{27}\)，这位侯爵表现得是个聪敏的学生。牛顿终于在 1693 年决定发表其流数的一个简略得近乎吝啬的概观（(XV), v. II, p.~391-396），这也无关紧要了；他的《原理》虽为一个多世纪提供了思考的食粮，但在无穷小这个领域，他已被赶上，并且在许多点上被超过。

然而，新体系的弱点至少在我们眼里是显而易见的。牛顿和莱布尼茨一举废除了两千年的传统，把首要的角色给予了微分，而把积分贬为仅仅是它的逆运算；需要整个\(XIX^{th}\)世纪、乃至\(XX^{th}\)世纪的一部分，才能通过把积分置于实变量函数一般理论及其现代推广的基础之上，重新确立一种公正的平衡（见《积分》一册的历史注记）。视角上的这种颠倒，也造成了不定积分以牺牲定积分为代价而占据的过分、几乎排他的地位，这一点在巴罗那里已可见到，在牛顿和莱布尼茨那里尤为明显：在这方面，同样须由\(XIX^{th}\)世纪来把事情摆正。最后，莱布尼茨特有的对符号作形式操作的倾向，在整个\(XVIII^{th}\)世纪继续变本加厉，远远超出了当时分析的手段所能辩护的限度。特别应当承认，莱布尼茨的微分概念说实话并没有意义；\(XIX^{th}\)世纪初它陷入了声名扫地之地，只是慢慢地才恢复名誉；一阶微分的运用虽已完全合法化，但高阶微分尽管十分有用，至今还没有真正得到正名。

不管怎样，微分学和积分学的历史自 17\(^{th}\) 世纪末起分为两个部分。其一涉及这门演算的应用，它们始终越来越丰富、越来越多、越来越多样。在前面谈到的平面曲线微分几何、微分方程、幂级数、变分学之外，又添上了挠曲线的微分几何，随后是曲面的微分几何、多重积分、偏微分方程、三角级数、众多特殊函数的研究，以及许多其他类型的问题，它们的历史将在专论它们的各册中阐述。我们在此只处理那些有助于打磨、深化和巩固无穷小演算基本原理本身的工作，就其涉及实变量函数的方面而言。

从这个观点看，18\(^{th}\) 世纪中叶的几大部头论著提供的新东西不多。苏格兰的麦克劳林\(^{28}\)和欧洲大陆的欧拉（XXV a 和 b）都忠实于各自所承的传统。的确，前者努力对牛顿的概念稍加澄清\(^{29}\)，而后者把莱布尼茨的形式主义推到极致，却同莱布尼茨和泰勒一样，满足于让微分演算建立在一个由差分演算出发的非常含混的极限过程之上，不过他对差分演算本身作了非常细致的阐述。但尤其重要的是，欧拉完成了莱布尼茨的事业：引入了 e、i 和三角函数沿用至今的记号，并推广了记号\(\pi\)。另一方面，尽管他多半不在函数与解析式之间作区分，但在讨论三角级数和振动弦问题时，他坚持不能局限于这样定义的函数（他称之为“连续的”），必要时也须考虑任意的、即“不连续的”函数，它们由一段或几段曲线弧经验地给出（(XXV c), p.~74-91）。最后，尽管这稍微超出了我们的框架，这里不可能不提到他把指数函数推广到复平面的工作，由此他得到了联系指数函数与三角函数的著名公式，以及复数对数的定义；对数与反圆函数之间那个著名的类比——用 17\(^{th}\) 世纪的语言说，即圆与双曲线的求积之间的类比——在那里得到了最终的阐明；这个类比格雷古瓦·德·圣樊尚已经观察到，惠更斯、尤其是格雷戈里使它精确化，而在莱布尼茨和伯努利那里，它出现在\(\frac{1}{1 + x^{2}} = \frac{i}{2(x + i)} - \frac{i}{2(x - i)}\)的形式积分之中。
% semantic-fragments: legacy-input-seam
\relax{}
与此同时，达朗贝尔——无论在数学中还是在其它领域，他都敌视一切玄奥——在一些出色的文章（XXVI）中以极为清晰的方式定义了极限和导数的概念，并竭力主张，归根结底这就是无穷小演算的全部“形而上学”。但这一明智的忠告并未立即产生效果。拉格朗日的宏篇巨著（XXVII）代表着一种尝试：把分析学建立在最有可议之处的牛顿派概念之一，即把任意函数的概念与可展开为幂级数的函数的概念混为一谈的那个概念之上，并由之导出微分的概念（通过考察级数中一阶项的系数）。当然，像拉格朗日这样水准的数学家在此不可能不获得重要而有用的成果，例如（以一种实际上与刚才指出的出发点无关的方式）给出带积分形式余项的泰勒公式的一般证明，并借助中值定理对它作出估计；此外，拉格朗日的工作还是魏尔斯特拉斯在复变函数论中的方法的源头，也是形式级数的现代代数理论的源头。然而，就其直接目的而言，它与其说是进步，不如说是退步。

相反，随着柯西的教学著作（XXVIII），人们重新站到了坚实的地面上。他实质上像我们今天这样定义了函数，尽管所用语言仍略嫌含糊。极限概念被一劳永逸地固定下来并被取为出发点；连续函数（现代意义下的）与导数的概念以及它们的主要初等性质都由之立即导出；而导数的存在性不再是一条信条，而是变成一个可用分析学的寻常方法来研究的问题。说实在的，柯西本人对此几乎未曾关注；另一方面，尽管波尔查诺从同样的原理出发，构造了一个在任何一点都没有有限导数的连续函数的例子（XXIX），这个例子并未发表，问题直到魏尔斯特拉斯 1872 年的著作（以及他 1861 年的课程）才得到公开解决（XXXII）。

至于积分学，柯西的工作意味着回到古代以及 XVII \(^{th}\) 世纪的健全传统，但所依赖的技术手段仍不完善。定积分在很长时期内被贬居次要地位，此时重新成为最根本的概念；对于它，柯西最终采用了傅里叶提出的记号 \(\int_{a}^{b} f(x) dx\)（以代替欧拉有时使用的累赘记号 \(\int f(x) dx \left[ \begin{array}{c} x = b \\ x = a \end{array} \right]\)）；而为了定义它，柯西回到了穷竭法，或者用我们的说法，回到“黎曼和”（把它称作阿基米德和或欧多克索斯和或许更好）。诚然，XVII \(^{th}\) 世纪从不认为应当把面积概念——在他们看来，这概念至少与不可公度的实数概念一样清楚——交给批判性的审查；但是“黎曼”和向曲线（哪怕只是单调曲线或分段单调曲线）下方面积的收敛，对于 XVII \(^{th}\) 世纪所有关心严格性的作者，如费马、帕斯卡、巴罗而言，都是熟悉的思想；而 J. 格雷戈里凭借他对取极限问题的思考以及他对已经十分抽象的“区间套”原理形式的熟谙而准备得格外充分，据推测他甚至已草成一份细致的证明，它一直未获发表（（XVII bis），第 445-446 页），倘若柯西知道有此证明，它几乎可以不加修改地为柯西所用\(^{30}\)。可惜对他而言，柯西竟想对任意的连续函数证明积分的存在性，即“黎曼和”的收敛性；他的证明如果以闭区间上连续函数一致连续的定理为基础，本是对的，但由于缺少这一概念，它不具任何证明效力。狄利克雷在撰写他那关于三角级数的著名论文时似乎也未察觉这一困难，因为他把所论的定理引作“容易证明”（（XXX），第 136 页）；诚然，归根结底他只把它用于有界的分段单调函数；黎曼更为审慎，在需要使用他关于“黎曼和”收敛的必要充分条件时，只提到这类函数（（XXXI），第 227-271 页）。在海涅建立起一致连续性定理（参看《一般拓扑学》，II，历史注记，第 217 页）之后，这个问题当然不再有任何困难；达布在 1875 年他关于不连续函数积分的论文（XXXIII）中便轻易地解决了它，在那篇论文里，他在许多点上恰与大约同时问世的 P. 杜·布瓦–雷蒙的重要研究不谋而合。这样一来，连续函数积分的线性便第一次得到了证明，而且这次是决定性地证明了。另一方面，由塞德尔等人在 1848 年引入、又被魏尔斯特拉斯善加利用的序列或级数的一致收敛概念（参看《一般拓扑学》，X，历史注记，第 347 页），已使人们能够为级数的逐项积分以及积分号 \(\int\) 下的求导奠定一个坚实的基础——诚然，所设条件稍嫌苛刻——同时等待着我们在此无须谈及的现代理论，它们以暂时最终的方式澄清了这些问题。

这样我们就到达了古典无穷小演算的最后阶段，即 XIX \(^{th}\) 世纪末那些伟大的《分析教程》所代表的阶段；按我们的观点，其中若尔当的教程（XXXIV）占有最崇高的位置，部分是出于审美上的理由，但也因为它虽是对古典分析成果的出色陈述，却在多方面预示了现代分析并为它准备了道路。若尔当之后到来的是勒贝格，那就进入本书另一册的主题了。

